\documentclass{exam}

\usepackage[margin=1in]{geometry}

\usepackage{comment}
\usepackage{graphicx}
\usepackage{listings}

\pagestyle{headandfoot}
\firstpageheader{CSCI 221}{}{C++ Intro Assignment}
\runningheader{CSCI 221}{}{C++ Intro Assignment}


\begin{document}
\title{C++ Intro Assignment (Learning)}
\author{CSCI 221: Computer Science Fundamentals 2}
\date{Spring 2026}
\maketitle


This assignment is an opportunity to test your understanding of C++ and receive feedback. 
Point values are assigned so that you can differentiate between large and small mistakes, but this assignment does not affect your grade. 
To submit your assignment, please submit your code, makefile, and writeup. 
For full credit, your code should check for invalid input. 

\textbf{Due Date:}
Monday, April 20th at 1:00 pm. 

\begin{questions}

\question[5]
\textbf{Hello World.}
\begin{parts}
\part[3]
Just like in C, your first step is to write and compile a basic working program. 
We will start with the Hello World program we used in the first lab. 
The code is below as a reference. 
\begin{verbatim}
#include <stdio.h>

int main (int argc, char *argv[]) {
  printf(``Hello, World!'');
  return 0;
}
\end{verbatim}
Get this code to compile using a C++ compiler. 
Make sure that your code is in a .cpp file rather than a .c file. 

\part[2]
Modify your code such that it uses a C++ style print statement rather than a C style print statement. 
Ensure your resulting code is correct. 
\end{parts}

%\question[]
%\textbf{Booleans.}

\question[10]
\textbf{Pointers.}
For the following, make sure to use C++ style rather than C style. 
\begin{parts}
\part[2]
Create a pointer to an integer. 
Don't initialize it. 
Write code that checks the value of that pointer correctly. 

\part[2]
Write code that allocates an integer and sets the pointer to point to that integer. 

\part[2]
Write code that deallocates the integer you just allocated. 

\part[2]
Write code that allocates an array of integers of length 17 and sets the pointer to point to that array. 

\part[2]
Write code that deallocates the integer array you just allocated. 

\end{parts}

\question[6]
\textbf{Functions.}
\begin{parts}
\part[2]
Write a function that takes a pointer to an integer and an integer and increments the value pointed to by the second value. 

\part[2]
Modify your function from the previous part to use C++ pass by reference semantics rather than a pointer. 

\part[2]
Modify your function so that if no value is passed in the second argument, it increments the first argument by one. 

\end{parts}

\question[5]
\textbf{Namespaces.}
\begin{parts}
\part[3]
Create a function \texttt{Update} inside the \texttt{Grading221} namespace that takes as input two doubles and returns a double. 
The returned value should be the first value if the first is larger, or the average if the second is larger. 

\part[2]
Create two different variables named \texttt{band}. 
One is used to refer to a ring-like object. 
One is used to refer to a musical group. 
Their types are unimportant, but your code should be able to access both variables inside the same scope. 
To show this, create a function that uses both variables. 

\end{parts}

\question[6]
\textbf{Strings and IO.}
\begin{parts}
\part[2]
Write a function that asks the user for an integer value, increments it by one, then prints out the result. 

\part[2]
Modify the function from the previous part to take a string from \texttt{cin}, then explicitly convert it to an integer using \texttt{stoi}. 
Similarly, explicitly convert the output from an integer to a string using \texttt{to\_string} before outputting it. 

\part[2]
Create another function that performs the same process, but this time uses the older string type conversion. 
That is, explicitly convert the user input from a string to an integer using a \texttt{stringstream}.
You will need to create a \texttt{stringstream} from the string, and then read out the integer like it was from \texttt{cin}. 
Similarly, convert your output from an integer to a string using a \texttt{stringstream}. 
You will need to have an empty \texttt{stringstream} read in your output and the call the \texttt{.str} function on the stream. 

\end{parts}

\question[28]
\textbf{Objects.}
\begin{parts}
\part[2]
Create a \texttt{CharacterCount} struct that contains a string and an array of integers. 

\part[8]
Write a function that takes as input a string and a pointer to a \texttt{CharacterCount}. 
The function has no return value. 
For each character in the string of the \texttt{CharacterCount}, the function should increment the integer in the same position by the number of times that character appears in the input string. 
You can assume that the length of the array in the \texttt{CharacterCount} is the same as the length of the string in the \texttt{CharacterCount.} 
For example, if a \texttt{CharacterCount} contains the string \texttt{"abcrst "}, then its array should have length 7. 
We will assume the array contains all zeros. 
If the function is passed the \texttt{CharacterCount} above and the string \texttt{"test string"}, after completion the CharacterCount array should contain
\begin{verbatim}
[0, 0, 0, 1, 2, 3, 1]
\end{verbatim}
because the character \texttt{t} appears three times in the string, \texttt{r}
appears once, \texttt{s} appears twice, etc.

\part[2]
Modify your struct so that it is a class containing the same variables. 
Don't add any methods yet. 
Modify your function so that it works on the class. 

\part[6]
Modify your class so that its variables are private. 
Create getter and setter methods to enable outside code to access them. 
However, you should enforce the restriction of the data type on the setter methods. 
In particular, when the string changes, you should allocate a new array of the appropriate length containing all zeros. 
In the setter for the array, you should make sure that you only accept arrays of the appropriate length. 
You should also make a function to set individual array values. 
Consider overriding the [] operator to do this. 
Modify your function to use these new methods. 

\part[2]
Create a constructor for your class that initializes all of the variables based on arguments containing the value to set them to. 

\part[2]
Create a default constructor for your class. 

\part[2]
Create a copy constructor for your class. 

\part[2]
Create a destructor for your class. 

\part[2]
If necessary, modify your code to ensure that all of your methods that should be constant are. 

\end{parts}

\question[10]
\textbf{Writeup.}
In addition to your code, you should submit a brief (~1-2 pages) writeup that describes the state of your code. 
You should discuss what is working, what is not, and any known errors. 
In this discussion, you should reference the tests you did and how they support your claims. 
You will need to provide a brief description of your tests for you to reference. 
Your writeup should be in a separate .txt or .pdf file. 

\end{questions}

\end{document}
